\documentclass{article}
\usepackage{amsthm}
\usepackage[dvipsnames]{xcolor}
\usepackage{avant}
\usepackage{fancyhdr}
\usepackage{tikz}
\usepackage{hyperref}
\usepackage[bottom]{footmisc}

\renewcommand{\familydefault}{\sfdefault}   % cambio font
\theoremstyle{definition}
\newtheorem*{definition}{Definizione}
\renewcommand{\figurename}{Figura}

\pagestyle{fancy}
\fancyfoot[L]{Basi di Dati}
\fancyfoot[C]{\thepage}
\fancyfoot[R]{Passarelli, Di Palma}

\setcounter{section}{-1}

\title{Reti e Laboratorio III}
\author{Angelo Passarelli, Giuseppe Di Palma}
\date{\today}

\begin{document}    
    \maketitle
    \begin{center}
        \includegraphics[scale=0.15]{Stemma_unipi.png}
    \end{center}
    \vspace{1cm}
    \begin{center}
        Appunti basati sulle lezioni e dispense delle professoresse Federica Paganelli \footnote{\url{http://pages.di.unipi.it/paganelli/}}
        e Laura Ricci \footnote{\url{https://pages.di.unipi.it/ricci/}}
    \end{center}
    \pagebreak
    \tableofcontents
    \pagebreak

    \begin{sloppypar}

        % SCRIVI QUI

    \end{sloppypar}

\end{document}  